% Manual de instalación y ejecución de PuellaGame.
\documentclass[11pt,letterpaper]{article}

\usepackage[utf8]{inputenc}
\usepackage[T1]{fontenc}
\usepackage[spanish,es-tabla]{babel}
\usepackage[margin=2.15cm,headheight=24pt]{geometry}
\usepackage{lmodern}
\usepackage{microtype}
\usepackage{graphicx}
\usepackage{booktabs}
\usepackage{tabularx}
\usepackage{array}
\usepackage{enumitem}
\usepackage{listings}
\usepackage[most]{tcolorbox}
\usepackage{fancyhdr}
\usepackage{titlesec}
\usepackage{xcolor}
\usepackage{tikz}
\usepackage{hyperref}

\definecolor{Navy}{HTML}{172A46}
\definecolor{Blue}{HTML}{246BCE}
\definecolor{Teal}{HTML}{0F9D8A}
\definecolor{Sky}{HTML}{EAF3FF}
\definecolor{Mint}{HTML}{E9F8F5}
\definecolor{Ink}{HTML}{263444}
\definecolor{Muted}{HTML}{607080}
\definecolor{Rule}{HTML}{D9E2EC}
\definecolor{CodeBg}{HTML}{F4F7FA}

\hypersetup{
  colorlinks=true,
  linkcolor=Blue,
  urlcolor=Blue,
  pdftitle={README PuellaGame - Equipo ÑIPGG},
  pdfauthor={Equipo ÑIPGG}
}

\setlength{\parindent}{0pt}
\setlength{\parskip}{5pt}
\setlist[itemize]{leftmargin=1.25em,itemsep=2pt,topsep=3pt}
\renewcommand{\arraystretch}{1.18}
\color{Ink}

\titleformat{\section}
  {\Large\bfseries\color{Navy}}{}{0pt}{}
  [\vspace{-0.35em}\color{Rule}\titlerule]
\titleformat{\subsection}
  {\large\bfseries\color{Blue}}{}{0pt}{}
\titlespacing*{\section}{0pt}{1.2em}{0.65em}
\titlespacing*{\subsection}{0pt}{0.9em}{0.35em}

\pagestyle{fancy}
\fancyhf{}
\lhead{\textcolor{Muted}{\small PuellaGame}}
\rhead{\textcolor{Muted}{\small Equipo ÑIPGG}}
\cfoot{\textcolor{Muted}{\small \thepage}}
\renewcommand{\headrulewidth}{0.4pt}
\renewcommand{\headrule}{\hbox to\headwidth{\color{Rule}\leaders\hrule height \headrulewidth\hfill}}

\lstdefinestyle{terminal}{
  basicstyle=\ttfamily\small,
  backgroundcolor=\color{CodeBg},
  frame=single,
  rulecolor=\color{Rule},
  framerule=0.5pt,
  xleftmargin=0.3em,
  framexleftmargin=0.6em,
  framexrightmargin=0.6em,
  framesep=6pt,
  breaklines=true,
  showstringspaces=false,
  columns=fullflexible
}

\newtcolorbox{summarybox}[1]{
  enhanced,
  colback=Sky,
  colframe=Blue,
  coltitle=white,
  fonttitle=\bfseries,
  title=#1,
  boxrule=0.7pt,
  arc=2mm,
  left=4mm,right=4mm,top=2mm,bottom=2mm
}

\newtcolorbox{notebox}[1]{
  enhanced,
  colback=Mint,
  colframe=Teal,
  coltitle=white,
  fonttitle=\bfseries,
  title=#1,
  boxrule=0.7pt,
  arc=2mm,
  left=4mm,right=4mm,top=2mm,bottom=2mm
}

\begin{document}
\pagenumbering{gobble}

\begin{titlepage}
  \thispagestyle{empty}
  \begin{tikzpicture}[remember picture,overlay]
    \fill[Navy] (current page.north west) rectangle ([yshift=-5.2cm]current page.north east);
    \fill[Blue] ([yshift=-5.2cm]current page.north west) rectangle ([yshift=-5.35cm]current page.north east);
  \end{tikzpicture}

  \vspace*{0.6cm}
  {\color{white}\large Universidad Nacional Autónoma de México\par}
  \vspace{0.2cm}
  {\color{white}\normalsize Facultad de Ciencias \textbullet{} Fundamentos de Bases de Datos\par}

  \vspace{2.6cm}
  {\color{Blue}\bfseries\large PRÁCTICA 02\par}
  \vspace{0.25cm}
  {\fontsize{30}{34}\selectfont\bfseries\color{Navy}PuellaGame\par}
  \vspace{0.18cm}
  {\Large\color{Muted}Manual de requisitos, compilación y ejecución\par}

  \vspace{0.8cm}
  \begin{summarybox}{Propósito del sistema}
    Prototipo de consola para administrar \textbf{sucursales, premios y clientes}.
    Ofrece operaciones de alta, consulta por identificador, modificación y
    eliminación, con persistencia en archivos CSV, validación de datos y manejo
    de excepciones.
  \end{summarybox}

  \vfill
  {\large\bfseries\color{Navy}Equipo ÑIPGG\par}
  \vspace{0.35cm}
  \begin{tabularx}{\textwidth}{>{\bfseries}p{0.7cm}X>{\raggedleft\arraybackslash}p{2.8cm}}
    \toprule
    No. & Integrante & Número de cuenta \\
    \midrule
    1 & Omar Alejandro Juárez Larrondo & 322245244 \\
    2 & Yahir León Bautista & 322176542 \\
    3 & Diego Hernández Gómez & 321069942 \\
    4 & Miguel Ángel Jiménez Ramírez & 119000887 \\
    5 & Ana Lilia Carballido Camacateco & 315314601 \\
    \bottomrule
  \end{tabularx}

  \vspace{0.7cm}
  {\small\color{Muted}Semestre 2027-I \hfill Versión 1.0\par}
\end{titlepage}
\pagenumbering{arabic}

\section{1. Requisitos previos}

\begin{tabular}{>{\bfseries\color{Navy}}p{3.6cm}>{\raggedright\arraybackslash}p{12.1cm}}
  \toprule
  Componente & Requisito \\
  \midrule
  Sistema operativo & Linux, macOS o Windows con terminal disponible. \\
  Java & JDK 25, requerido por \texttt{maven.compiler.release}. \\
  Maven & Apache Maven 3.8 o posterior. \\
  Codificación & UTF-8 para fuentes y archivos CSV. \\
  Conectividad & Necesaria en la primera compilación si Maven aún no tiene las dependencias. \\
  \bottomrule
\end{tabular}

Verificar las herramientas instaladas:

\begin{lstlisting}[style=terminal]
java --version
javac --version
mvn --version
\end{lstlisting}

\begin{notebox}{Directorio de trabajo}
Todos los comandos de este manual deben ejecutarse desde \texttt{Practica02/SRC}.
La aplicación busca los archivos persistentes en la ruta relativa
\texttt{datos/}; ejecutarla desde otro directorio impediría localizar los CSV
entregables.
\end{notebox}

\section{2. Organización del sistema}

\begin{tabular}{>{\ttfamily\color{Blue}}p{3.4cm}>{\raggedright\arraybackslash}p{12.3cm}}
  \toprule
  Ruta & Contenido \\
  \midrule
  model/ & Entidades \texttt{Sucursal}, \texttt{Premio} y \texttt{Cliente}. \\
  repository/ & Persistencia y conversión de registros CSV. \\
  service/ & Reglas de negocio y operaciones CRUD. \\
  ui/ & Menús, formularios y lectura validada de consola. \\
  validation/ & Validación reutilizable de entradas. \\
  datos/ & Archivos \texttt{sucursales.csv}, \texttt{premios.csv} y \texttt{clientes.csv}. \\
  test/java/... & Pruebas automatizadas de repositorios, servicios, interfaz y validación. \\
  Doc/apidocs/ & Documentación HTML generada con Javadoc. \\
  \bottomrule
\end{tabular}

\section{3. Compilación y verificación}

\subsection{Ejecutar todas las pruebas}

La forma recomendada de preparar el proyecto es limpiar los artefactos previos,
compilar las fuentes y ejecutar la suite completa:

\begin{lstlisting}[style=terminal]
cd Practica02/SRC
mvn clean test
\end{lstlisting}

Una ejecución correcta termina con \texttt{BUILD SUCCESS}. La versión entregada
incluye pruebas para las operaciones CRUD, persistencia CSV, validaciones y
navegación por consola.

\subsection{Compilar sin repetir las pruebas}

\begin{lstlisting}[style=terminal]
mvn compile
\end{lstlisting}

Las clases compiladas quedan en \texttt{target/classes}. El directorio
\texttt{target/} es regenerable y no debe editarse manualmente.

\subsection{Generar la documentación}

\begin{lstlisting}[style=terminal]
mvn javadoc:javadoc
\end{lstlisting}

El punto de entrada de la documentación es
\texttt{Doc/apidocs/index.html}. La configuración del proyecto trata las
advertencias de Javadoc como errores para evitar entregar documentación incompleta.

\section{4. Ejecución de PuellaGame}

Después de compilar, iniciar la aplicación con:

\begin{lstlisting}[style=terminal]
java -cp target/classes mx.unam.ciencias.nipgg.puellagame.app.Main
\end{lstlisting}

El menú principal permite seleccionar la entidad que se desea administrar:

\begin{lstlisting}[style=terminal]
=== PuellaGame ===
1. Gestionar sucursales
2. Gestionar premios
3. Gestionar clientes
0. Salir
\end{lstlisting}

\clearpage
\subsection{Operaciones disponibles}

Cada submenú ofrece las mismas operaciones:

\begin{tabular}{>{\bfseries\color{Navy}}p{2.7cm}>{\raggedright\arraybackslash}p{13cm}}
  \toprule
  Operación & Comportamiento \\
  \midrule
  Agregar & Solicita todos los campos, valida el identificador y rechaza duplicados. \\
  Consultar & Recibe la llave de la entidad y muestra todos sus datos. \\
  Editar & Conserva la llave indicada y reemplaza los demás campos validados. \\
  Eliminar & Solicita confirmación antes de retirar el registro del CSV. \\
  \bottomrule
\end{tabular}

Las opciones no válidas, valores no numéricos, enteros fuera de rango, correos
incorrectos y teléfonos que no tengan diez dígitos se informan sin cerrar el
programa. La captura se repite hasta recibir un valor aceptable.

\section{5. Persistencia y formato CSV}

Los cambios se escriben inmediatamente en un archivo independiente por entidad.
Cada archivo contiene una primera fila de encabezados y utiliza comas como
separadores:

\begin{tabular}{>{\ttfamily\color{Blue}}p{3.4cm}>{\raggedright\arraybackslash}p{12.3cm}}
  \toprule
  Archivo & Columnas \\
  \midrule
  sucursales.csv & idSucursal, nombre, direccion, telefono \\
  premios.csv & idPremio, nombre, descripcion, puntosRequeridos, existencias \\
  clientes.csv & idCliente, nombre, correo, telefono, puntosAcumulados \\
  \bottomrule
\end{tabular}

\begin{notebox}{Conservación de datos}
No se deben eliminar los encabezados ni editar los archivos mientras el programa
esté en ejecución. Los campos que contienen comas, comillas o saltos de línea se
escapan conforme al formato CSV antes de guardarse.
\end{notebox}

\clearpage
\section{6. Solución de problemas}

\begin{tabular}{>{\raggedright\arraybackslash\bfseries\color{Navy}}p{4.2cm}>{\raggedright\arraybackslash}p{11.5cm}}
  \toprule
  Mensaje o síntoma & Acción recomendada \\
  \midrule
  Maven no disponible & Instalar Maven y confirmar que esté disponible en \texttt{PATH}. \\
  Error de versión de Java & Seleccionar un JDK 25 y verificar \texttt{JAVA\_HOME}. \\
  No se encuentra \texttt{datos/} & Ejecutar el programa desde \texttt{Practica02/SRC}. \\
  ID duplicado & Elegir un entero positivo que no exista en el CSV correspondiente. \\
  Entidad inexistente & Confirmar la llave antes de consultar, editar o eliminar. \\
  CSV no legible & Revisar permisos, codificación UTF-8 y encabezados del archivo. \\
  \bottomrule
\end{tabular}

\section{7. Lista de verificación}

Antes de considerar preparada la parte de código de la entrega:

\begin{itemize}
  \item \texttt{mvn clean test} termina con \texttt{BUILD SUCCESS}.
  \item \texttt{mvn javadoc:javadoc} termina sin advertencias.
  \item El menú abre las tres entidades y permite las cuatro operaciones CRUD.
  \item Los tres CSV conservan sus encabezados y reflejan los cambios realizados.
  \item La aplicación informa errores de entrada sin finalizar inesperadamente.
  \item El directorio \texttt{Doc/apidocs/} contiene la documentación generada.
\end{itemize}

\vfill
\begin{center}
  \small\color{Muted}
  Equipo ÑIPGG \textbullet{} Práctica 02 \textbullet{} Fundamentos de Bases de Datos
\end{center}

\end{document}
